% 本文件由 LaTeX 排版 Agent 根据有效 translation.json 自动生成。
% author=Methodius; work_id=0627; title=西默盎与亚纳讲道词

\chapter{西默盎与亚纳讲道词}

% structure: p0001 source=h1 target=omitted reason=page-title-already-rendered-by-parent

在他们在圣殿相遇之日\par

虽然我之前在关于贞洁的对话中，已尽可能简要地为论童贞的讲道打下了充分的基础，然而今日的时节却将童贞之荣耀及其不朽荣冠的整个主题，呈献给教会养育的儿女们，作为愉快的默想主题。因为今日，神圣神谕的议会厅已敞开，预表这荣耀之日及其效果与结果的征兆，由神圣的宣讲者向聚集的教会宣读。今日，那古老而真实旨意的完成，在事实与行动上，荣耀地向世界彰显。今日，没有任何遮蔽，\BibleRef{格后 3:18}并以露着脸，我们如同在镜子中看见主的荣耀和神圣约柜本身的威严。今日，至圣的集会，承担着历代所期待的天上喜悦，将它传给人类。\BibleQuote{旧事已过} \BibleRef{格后 5:17}——新事物绽放如花，且永不凋谢。严苛的法律法令不再掌权，但主的恩宠统治着，借着救赎性的忍耐吸引众人归向祂。乌匝不再第二次\BibleRef{撒下 6:7}因胆敢触摸不可触摸之物而受无形惩罚；因为天主亲自邀请，谁还会怀着恐惧犹豫不前呢？祂说：\BibleQuote{凡劳苦和负重担的，你们都到我这里来。} \BibleRef{玛 11:28}那么，谁不奔向祂呢？不要让任何犹太人看着厄贝得东家的预像\BibleRef{撒下 6:10}主已\Quote{\Emph{明显来到自己的地方}。}祂坐在一个活生生的而非无生命的约柜上，如同坐在恩宠之座上，庄严地游行于大地。税吏触摸这约柜，离开时成了义人；娼妓接近它，仿佛被重塑，变得贞洁；癞病人触摸它，毫无痛苦地痊愈了。它不拒绝任何人，不回避任何人；它分施治愈的恩赐，自身却不沾染任何疾病；因为那喜爱并关心人的主，在其中安息。这些是这新恩宠的恩赐。这是那在太阳下发生的新奇事\BibleRef{德 1:10}——一件前所未有、后无来者的事。天主因祂对我们的怜悯所预定的，已经实现，祂因那如此适合祂的慈爱而使之成全。因此，神圣的号角有充分的理由吹响：\BibleQuote{旧事已过，请看，一切都成了新的。} \BibleRef{格后 5:17}我又当想什么，说什么才配得上这天呢？我正挣扎着接近那不可及者，因为对这位神圣贞女的追忆远远超越我所有的话语。因此，既然所需的颂词之伟大完全羞辱了我们有限的能力，就让我们投奔那未超出我们能力的赞美诗，并夸耀我们不可改变的失败，让我们加入基督羊群中欢庆节日的喜乐合唱。而你们，我神圣的听众们，请保持肃静，以便那满载真理的船只，通过耳朵的狭窄通道，如同进入理解的港湾，安然航行。我们守节，不是依照希腊神话的虚妄习俗；我们守的庆典不带任何荒谬或疯狂的众神宴饮，而是教导我们那超越万有的天主之可畏荣耀向人类的奇妙屈尊。\BibleRef{罗 9:5}\par

II. 来吧，依撒意亚，最庄严的宣讲者，最伟大的先知，智慧地向教会展开在荣耀中聚会的奥秘，激励我们卓越的宾客，使他们饱享持久的佳肴，以便你这位诚实的先知，将我们拥有的现实对照于你那面镜子，并为你的预言的结果而快乐地鼓掌。他说：\Quote{在乌齐雅王逝世那年，我看见主坐在崇高的宝座上，圣殿充满了祂的荣耀。\OriginalTerm{色辣芬}{Seraphim}侍立左右，各有六个翅膀。他们彼此呼喊说：\InnerQuote{圣哉，圣哉，圣哉，万军的上主！整个大地充满了祂的荣耀。}呼喊者的声音震动了门限，圣殿内充满了烟云。我说：我有祸了！我心中刺痛，因为我是个唇舌不洁的人，住在唇舌不洁的人民中间，因为我亲眼看见了君王——万军的上主。有一个色辣芬奉命到我这里来，手中拿着从祭台上用火钳取下的红炭，触及我的口说：\InnerQuote{看，这炭接触了你的口唇，你的罪孽已被消除，你的罪恶已得净化。}我又听见上主的声音说：\InnerQuote{我将派遣谁呢？谁肯为我们去呢？}我回答说：我在这里，请派遣我！祂说：\InnerQuote{你去对这人民说：你们听了，却不明白；看了，却不理解。}} 这些是先知借着圣神预先宣告的。亲爱的诸位，请思考这些话的力量。如此你们将理解这些圣事性象征的结果，并知道我们自己这次聚集是什么以及多么伟大。既然先知已预言了这奇迹，来吧，以最大的热情、欢欣和心灵的敏捷，并带着你们智力最敏锐的洞察力，一同前往著名的\Place{白冷}，在你们心中描绘一幅清晰分明的图像，将预言与实际事件的结果相比较。你们无需多言就能认识这事；只需注视那里正在发生的事。\Quote{一切智慧之事对于明达的人都是清楚的，对于获得知识的人都是正直的。}\BibleRef{箴 8:9} 因为，请看，如同一个因创造者的荣耀而崇高升起的宝座，童贞之母在那里预备好了，显然是给君王——万军的上主。在这宝座上，你们要看到主现在借着有罪的肉身来到你们中间。在这个童贞的宝座上，我再说，朝拜祂，祂现在借着这条新的、永远可敬的道路来到你们中间。用信德之眼环顾四周，你们会发现在祂周围，按照他们的行列秩序，有如王家和司祭的色辣芬队伍。这些作为祂的卫兵，常习惯地陪伴他们君王的临在。因此，在这里，他们不仅被说成以赞美歌颂那神圣一体的神圣本体，而且也被说成赞美那神圣三位一体中一位的全可敬拜的荣耀，这位现在因天主在肉身显现而降临于大地。他们说：\Quote{整个大地充满了祂的荣耀。} 因为我们认为，与为了我们而成为人的圣子一起，按照祂旨意的善意，圣父——祂在神圣本性上与圣子不可分离——也在场，同质的圣神也在场。因为正如神圣神谕的解释者\Person{保禄}所说：\Quote{天主在基督内使世界与自己和好，不再追究他们的过犯。}\BibleRef{格后 5:19} 他以此表明圣父在圣子内，因为同一个旨意在祂们内运作。\par

III. 因此，噢，这庆节的爱好者，当你仔细思量白冷的荣奥秘迹（这些是为你的缘故而成就的）时，你要欣然加入那正在为你庄严庆祝救恩的天军。\BibleRef{撒下 6:14} 正如从前达味在约柜前所做，你也要在这童贞的宝座前欢欣舞蹈。用欢快的歌声咏唱那永远无处不在的主，以及那从特曼来的祂，\BibleRef{哈 3:3} 如先知所说，乐意显现于人类，且是在肉身中。同梅瑟一起说：\BibleQuote{祂是我的天主，我要光荣祂；我父亲的天主，我要颂扬祂。} \BibleRef{出 15:2} 然后，在你感恩的赞美诗之后，我们要有益地探究是什么原因激起荣耀的君王在白冷显现。祂对我们的怜悯迫使那位不可被强迫者，在白冷以一个人类身体出生。但有什么必要使祂，一个吃奶的婴孩，那虽在时间中却不被时间所限，那虽裹在襁褓中却不被它们束缚，有什么必要使祂成为流亡者、离开祂的祖国？如果你确实想要知道这个，你这至圣的会众，天主圣神已吹息在你们身上，请听梅瑟清楚地向人民宣告，仿佛激励他们认识这不寻常的诞生，说：\BibleQuote{凡开胎的男性，应称为圣于主。} \BibleRef{出 31:19} 噢奇妙的境况！\BibleQuote{噢，天主智慧和知识的富藏何其深奥！} \BibleRef{罗 11:33} 诚然，法律和先知的主按照祂自己的法律行一切事，并不使法律无效，而是成全它，并更将祂恩宠的开始与法律的成全联系起来是合适的。因此，那高于法律的母亲服从了法律。而这位圣洁无玷者，遵守为不洁者所定的四十天期限。那使我们从法律中获得自由者，反而成了法律的属下；且为那圣化我们者，献上了一对洁净的鸟，\BibleRef{路 11:24} 为证明那些近前来的人是洁净无瑕的。既然那次分娩是污秽的，且不需要赎罪的牺牲，依撒意亚是我们的证人，他向普天之下清楚宣告：\BibleQuote{她未受产痛，} 他说，\BibleQuote{在疼痛来到之前，她已生产，她逃脱了，生了一个男孩。} \BibleRef{依 66:7} 谁曾听过这样的事？谁曾见过这样的事？所以，那至圣的童贞母亲，甚至在分娩之前就完全免除了妇女的惯例：无疑，为使圣神聘娶她归于自己，并圣化她，使她能在与男人无交合的情况下怀孕。她生下了她的长子，即天主的独生子，我说，那在天上毫无母亲地从圣父的实体中发出作为独生子，并保守祂本性合一的童贞完整无损、不可分割；那在地上，在童贞女的婚房中，像新郎一样，以不可分离的结合将亚当的本性归于自己，并保守祂母亲的无玷纯洁未受损伤——简言之，那在天上无玷而生，在地上以完全不可言说方式所生者。但言归正传。\par

四、因此，先知将童贞玛利亚从纳匝肋带来，使她能在白冷生下她带来救恩的婴孩，又将她带回纳匝肋，为向世界彰显生命的希望。故此，天主的约柜从白冷的客栈起程，因为在那里，祂为法律偿还了四十天的债，这债不是出于正义，而是出于恩宠；然后，她安息在熙雍的山上，将她纯洁的怀中如同崇高的宝座、超越人性的宝座，接受万有的君王，并将祂呈献给天主圣父，作为祂宝座的同享者和祂本性不可分离的一位，连同祂从她身上所取的纯洁无瑕的肉体。圣母亲自登上圣殿，为向法律展示一个新奇而奇妙的景象：那长久期待的婴孩，祂开启了童贞女的子宫，却未破坏童贞的界限；那超越法律的婴孩，却履行了法律；那既在法律之前又在其后的婴孩；简言之，那超越自然律而由她降生成人的婴孩。因为在其他情况下，每个子宫首先通过与男人的结合而开启，并由他的种子受孕，获得受孕的开端，然后通过完成分娩的阵痛，最终按照创造主天主的上智安排，生出有理性、与本性和谐的后裔。因为天主说：\BibleQuote{你们要生育繁殖，充满大地。}而这位童贞女的子宫，未经开启，也未受孕，却生出了一个超越本性而又与之同源的子嗣，无损于不可分割的统一性，以致奇迹更令人惊叹，童贞的特权也完整无损。因此，她——那位比圣殿更崇高的——登上圣殿，身披双重荣耀：即无玷童贞的荣耀和不可言喻的生育之福，法律的祝福和恩宠的圣化。因此，目睹这情景的先知说：\BibleQuote{整个房屋充满了祂的光荣，撒拉弗侍立在祂周围；他们彼此呼喊说：圣！圣！圣！万军的上主！祂的光荣充满大地！}\BibleRef{依 6:3}同样，蒙福的先知哈巴谷也优美地歌唱说：\BibleQuote{在两只活物中间，你将被人认识；当岁月临近时，你将被认出；当时间来到时，你将显现。}请看，我恳求你，圣神的话是何等精确。祂谈到认识、认出、显现。关于第一点：\BibleQuote{在两只活物中间，你将被人认识}\BibleRef{哈 3:2}，这指的是在天主法律时期，天主的光荣曾停留在至圣所内约柜的盖上，在预象的革鲁宾之间，正如祂对梅瑟说：\BibleQuote{在那里我将向你显现。}\BibleRef{出 25:22}这也指天使的会聚，现在因救主亲自在天主性上永远可崇拜的肉体显现而前来迎接我们，尽管祂的本性是我们不能看见的，正如依撒意亚先前已宣告的。但祂说：\BibleQuote{当岁月临近时，你将被认出}，意思如先前所说，是指在肉身内的天主、我们的救主，那否则对凡眼不可见的荣耀认识；正如伟大的神圣奥秘诠释者保禄在某处说：\BibleQuote{但时期一满，天主就派遣了祂的儿子，生于女人，生于法律之下，为把在法律之下的人赎出来，使我们获得义子的地位。}\BibleRef{迦 4:4}至于随后的话：\BibleQuote{当时间来到时，你将显现}，如果人仔细将心目光投向我们现在庆祝的庆节，这还需要什么解释呢？\Quote{因为那时你要显现，}祂说，\Quote{如同在君王的坐骑上，由你纯洁无瑕的母亲怀抱，在圣殿中，在你所取肉身的恩宠和美丽中。}先知为了更清晰而总结这一切，简短地呼喊说：\BibleQuote{上主在祂的圣殿内；}\BibleRef{哈 2:20}\BibleQuote{全地都当在祂面前战栗。}\par

五．与妳相关的奥迹确实可畏，童贞之母，你这属神的宝座，受光荣并被天主视为相称。你在天上和世人的眼前，产生了一个卓越的奇事。这是此事的明证，一个不可辩驳的论证：在你超自然生育的新奇之时，天使在地上歌唱：\BibleQuote{天主受享光荣于高天，主爱的人在世享平安} \BibleRef{路 2:14}，他们以三重歌声引入三重圣德。你在妇女中是有福的，你这至受天主祝福者，因为藉着妳，大地充满了天主的那个神圣光荣；正如在圣咏中所唱的：\BibleQuote{愿上主天主，以色列的天主受赞美；愿他的光荣充满大地！阿们。阿们。} 而门限的基石，先知说，因那呼喊者的声音震动，这所指的是约柜前的帐幔，它预表了妳，为使我认识真理，也为了藉着之前的预象和预表，让我学会以敬畏和战兢前来尊敬与你相关的神圣奥迹；并且藉着法律的这先前的阴影描绘，使我得以约束自己，不敢大胆无礼地用定睛的目光注视那在其不可理解性中远远高居万有之上的那位。因为如果对于约柜——那是妳圣洁的肖像和预象——天主曾给予如此尊荣，以致除司祭阶层外，无人可以接近或获准瞻仰它，并有帐幔隔开，将前厅如王后般保留，那么，对于我们这些受造物中最微末的人，理应向妳——妳确实是王后——向妳，天主的活约柜，立法者——向妳，那容纳不可容纳者的诸天——呈上何等、何等的敬意？因为，圣童贞，妳已如光明日晨照世界，并生出了义德的太阳，那可憎的黑暗恐怖已被驱散；暴君的权力已被摧毁，死亡已被消灭，地狱已被吞没，一切敌意已在和平面前消散；有害的疾病如今退去，因为救恩已显现；整个宇宙已被纯洁而清晰的真理之光充满。这些事在\BibleBook{}{雅歌}中撒罗满有所暗示，并如此开始：\BibleQuote{我的爱者属于我，我属于我的爱者；他在百合花间牧羊，直到天破晓，阴影遁去。} \BibleRef{歌 2:16} 既然天主之神已在熙雍显现，他荣美的光辉已在耶路撒冷出现，\BibleQuote{光明为义人升起，为心地正直的人预备了喜乐。} 依照有福的达味，成全者的成全者及主，藉着圣神，召叫了法律的教师和仆役，为服事并证言所成就的事。\par

VI. 因此，年迈的西默盎脱下肉身的软弱，穿上希望的力量，急忙迎接法律的执事，拥有权柄的导师，亚巴郎的天主，依撒格的保护者，以色列的圣者，梅瑟的导师；我说，就是祂，曾应许向他显示祂神圣的降生成人，如同祂的背面；\BibleRef{出 33:23}祂在贫穷中却是富有的；祂在婴孩时期已先于万世；祂虽被看见，却是不可见的；祂虽可理解，却是不可理解的；祂虽微小，却超越一切宏大——既在圣殿内，又在至高之天——在君王宝座上，又乘坐革鲁宾的战车；祂既在上也在地下，连绵不断；祂既有仆人的形像，又有天主圣父的形像；祂是臣仆，却又是万王之王。他完全沉浸在渴望、希望和喜乐中；他不再属于自己，而是属于那位他所期待者。圣神已向他宣报了喜讯，在他抵达圣殿之前，他藉着理智的眼目被高举，仿佛已拥有了他所渴求的，他欢欣踊跃。就这样被引领，他匆忙中脚步踏着空气，来到那历来被视为神圣的圣所；但他并未注目圣殿，却向圣殿的主宰伸出他圣洁的双臂，唱出适合这喜庆场合的诗歌：\Quote{我切慕你，上主，我祖的天主，仁慈之主，你屈尊就卑，出于你的光荣和良善，你供给一切，出于你恩宠的俯就，你借此倾向我们，作为带来和平的中保，建立天地间的和谐。我寻求你，伟大的万有创造者。我切望期待你，你用你的圣言包容万有。我等候你，生命与死亡的主宰。我仰望你，法律的赐予者，法律的继承者。我饥渴于你，你使死者复生；我口渴于你，你使困乏者得安慰；我渴慕你，世界的创造者与救赎者。\BibleRef{依 43:10}你是我们的天主，我们朝拜你；你是我们的圣殿，我们在你内祈祷；你是我们的立法者，我们服从你；你是万物的元始天主。在你之前，没有别的神由天主圣父所生；在你之后，也没有别的子，与圣父同性同体，同享光荣。认识你就是完美的正义，认识你的权能就是不朽的根源。\BibleRef{智 15:3}你就是那一位，为了我们的救恩，成了宝贵的、可敬的角石，预先向熙雍宣告了。因为万有都置于你之下，作为他们的因由和创造者，你从无中创造了万有，并使不稳的得以坚定结合；你是受造物的联系和保护者；你是不同本性事物的塑造者；你用智慧和稳定的手掌握宇宙的舵；你是完美秩序的准则；你是和谐与平安不可辩驳的纽带。因为我们在你内生活、行动和存在。\BibleRef{宗 18:28}因此，上主我的天主，我要光荣你，我要赞美你的名；因为你行了奇妙的事；你古时的计谋是信实和真理；你以威严和尊荣为衣。因为对君王来说，有什么比绣花紫袍和闪亮的王冠更华美呢？对喜爱人的天主来说，有什么比这仁慈的人性摄取更为宏伟呢？它用灿烂的光芒照亮了坐在黑暗和死影中的人。\BibleRef{依 42:7}那位暂时的君王，你的仆人，曾恰当地歌颂你为永恒之王，说：你比世人更美，你在人中是真实的天主和人。因为你藉着降生成人，以正义束腰，以信实膏抹你的脉络，你自己就是正义和真理，是万有的喜乐与欢腾。\BibleRef{依 11:5}因此，诸天哪，今天与我一同喜乐，因为上主向祂的百姓显示了仁慈。是的，愿云彩将正义的甘露降在世界之上；愿大地的根基向阴府吹响号角，因为沉睡者的复活已来临。\BibleRef{依 45:8}也愿大地使怜悯向它的居民滋生；因为我满得安慰，我极其喜乐，因为我已看见了你，人类的救主。\BibleRef{格后 7:4}}\par

VII. 当那老者如此欢欣鼓舞，怀着极其伟大而神圣的喜乐时，先知\Person{依撒意亚}先前以预象所说的事，如今由天主之母明显地实现了。因为她从自己纯净无瑕的祭台上，拿起那活生生的、不可言喻的、披上了人肉的煤炭，以她神圣的双手的怀抱——如同火钳一般——将祂递给那位义人，并似乎以这样的话劝勉他说：\Quote{请接收，可敬的长者，您，司祭中最卓越者，接收主吧，并尽享您那不再孤寂无望的希望。请接收，人中至为杰出者，那永不枯竭的宝藏，那些不能被夺去的财富。请抱紧，人中至为智慧者，那不可言喻的权能，那无法探究的能力，唯有祂能支撑您。请拥抱，圣殿的侍奉者，那无限的伟大，那无与伦比的力量。请您环绕着生命本身，并生活，至为可敬者。请紧贴不朽，并得以更新，至为正义者。此举并不孟浪，因此您不要退缩，至为圣洁者。请以您所渴望的祂来满足自己，并在那位被赐予——更确切地说，是那位将祂自己赐予您的那一位中欢欣，至为神圣者。请喜乐地汲取您的光明，至为虔诚者，从正义的太阳那里，祂通过无瑕的肉身这面镜子向您照耀。不要害怕祂的温和，也不要让祂的仁慈令您恐惧，至为蒙福者。不要害怕祂的宽厚，也不要因祂的良善而退缩，至为谦逊者。请欢欣地与祂结合，不要迟延服从祂。向您所说和所展示的，并无鲁莽之处。因此，请不要犹豫，至为文雅者。我主恩宠的火焰并不吞噬，而是光照您，至为正义者。让那以我作为预象的荆棘，使您明白这一点，您这位在法律上受教最深者——那荆棘象征着那尚无实质的火焰的真实。\BibleRef{出 3:2} 让那如同散发露珠的微风一般的火窑，说服您，老师，关于这一奥秘的救恩计划。\BibleRef{达 3:21} 此外，让我的子宫为您作证，其中容纳了那从未被容纳于任何事物中的那一位，这证明圣言降生成人所取的人性实体。如今号角的响声不再惊吓近前来的人，也不再有冒烟的山使接近者恐惧，法律也不再无情地惩罚那些敢于触摸的人。\BibleRef{出 19:16} 眼前的事物诉说着对人的爱；显现的事物诉说着天主的屈尊就卑。因此，请感恩地接收来到您面前的天主，因为祂将除去您的过犯，彻底净化您的罪。在您身上，让世界的洁净首先，如同在预象中，得以实现。在您身上，并藉着您，让那因恩宠而成的成义预先显示给外邦人。您堪当这使人复活的初果。您善用了法律。今后，请运用恩宠。您已因文字而疲惫；在圣神中更新吧。脱去那旧的，穿上那新的。因为关于这些事，我想您并非不知道。}\par

八、因此，那位义人胆大起来，顺从天主之母（她是在关涉人类的事上作天主的婢女）的劝勉，用他年迈的双手接抱了那位在婴孩中仍是亘古常在者，并祝福了天主，说道：\BibleQuote{上主，现在可照祢的话，让祢的仆人平安离去；因为我亲眼看见了祢的救恩，就是祢在万民面前所预备的：为作启示外邦人的光明，和祢百姓以色列的荣耀。}\BibleRef{路 2:29}我从祢领受了毫无掺杂痛苦的喜乐。欢欣地接纳我吧，上主，我歌颂祢的仁慈与怜悯。祢将这心灵的喜乐赐给了我。我以喜乐向祢献上感恩的祭。我认识了天主之爱的力量。因为，为了我，由祢所生、不可言喻、毫无玷污的天主成了人。我认识了祢对我们之爱与关怀的难以言喻的伟大，因为祢派遣了祢自己的心肠来解救我们。如今，我终于明白了从撒罗满所学到的：\BibleQuote{爱情如死亡一样坚强：因为藉着它，死亡的毒刺将被除灭，藉着它，死人将得见生命，藉着它，甚至连死亡也将学到什么是死亡，并被迫停止它对我们所行使的统治。藉着它，那作我们邪恶之源的蛇也将被擒获并淹没。}\BibleRef{歌 8:6}祢使我们认识了祢的救恩，上主，为我们生出了和平的植物，我们不再误入迷途。祢使我们知道，上主，祢并未至终忽略祢的仆人；祢，良善者，也未曾完全忘记祢双手的工程。因为祢怜悯我们卑微的境况，将祢那无穷尽、与祢本性相同的良善丰沛地倾注在我们身上，藉着祢的独生子救赎了我们，祂与祢毫无变化地相似，并与祢同体；祢认为将拯救和造福祢仆人的工作托付给一个仆人，或使那些冒犯者藉着一个中间人与祢和好，是与祢的威严和良善不相称的。但是，藉着那与祢同体的光明，祢光照了那些坐在黑暗和死影中的人，好使他们在祢的光中得见知识的光明；并且祢认为藉着我们的主和创造者，重新塑造我们归于不死是好的；祢仁慈地赐给了我们返回乐园之路，藉着那将我们从乐园的喜乐中分离出来的那位；并且藉着那位拥有赦罪权柄的，祢\BibleRef{谷 2:10}涂抹了那反对我们的债卷。\BibleRef{哥 2:4}最后，藉着那位分享祢宝座、不能与祢的神性分离的那位，祢赐给了我们和好的恩赐和坦然无惧亲近祢的权利，好使藉着那不认识任何主权的主，藉着真实而全能的天主，这么多如此伟大祝福的签署认可（可以说是）构成恩宠的成义恩赐，成为那些获得怜悯者确定无疑的权利。正是这件事，先知早已预言说：没有使节，也没有天使，乃是主亲自拯救了他们；因为祂爱他们，顾惜他们，并扶起他们，高举他们。这一切，并不是由于我们所行的义行\BibleRef{铎 3:5}，也不是因为我们爱了祢——因为我们第一个尘世的先祖，在乐园愉悦的居所中受尊荣，却藐视了祢神圣而救人的诫命，被判定为不配那生命之地，并将他的种子与罪的私生子枝条混杂，使之极其软弱——但是，祢，上主，出于祢自己，并出于祢对祢亲手所造的受造物不可言喻的爱，坚定了祢对我们的怜悯，并怜悯我们与祢的疏远，看到我们堕落\BibleRef{若 4:9}，祢就动心，将我们纳入怜悯之中。因此，今后为我们亚当族类确立了一个喜乐的庆节，因为亚当的第一创造者出于自己的自由意志成了第二亚当。主我们天主的光辉降临与我们同住，使我们得见天主的面，并得救。因此，上主，我求祢准我离去。我已看见了祢的救恩；让我从我文字的轭中得释放。我已看见了永恒的王，没有后继者；让我从这奴役而沉重的锁链中得自由。我已看见了那本性的上主和拯救者；那么，愿我能获得祂对我的释放令。使我从定罪的轭下释放，将我置于成义的轭下。救我脱离诅咒的轭和那叫人死的文字；\BibleRef{格后 3:6}并将我登记在那些有福的行列中，他们藉着祢这位真实之子（祂与祢同荣同能）的恩宠，已被收纳为义子。\par

九. 那么，他说，让我到目前为止简要说的这些，作为我向天主的感恩之献，便足够了。但我要对你说什么呢，啊，母亲童贞女与童贞母亲？因为赞美她——这非人所造者，超乎人力之所能。因此，我要以闪耀在你周围的诸天使恩赐的光辉，照亮我贫乏的昏暗，并从不朽的草地上采撷花冠，奉献给你，献上属于你自己的，为你神圣且受天主加冕的头顶。我要以你祖传的颂歌迎接你，啊，\Person{达味}之女，以及\Person{达味}的主和天主之母。因为用属于他人的东西来装饰你——你在自己的荣耀中已然超卓——是既卑劣又不祥的。因此，啊，最仁慈的圣母，请接受珍贵的礼物，那些唯独适合你的礼物。啊，你被高举于万代之上，在一切受造之物——有形与无形的——之中，闪耀为最可尊敬者。\Person{叶瑟}的根是蒙福的，\Person{达味}的家是三次蒙福的，你从中生长出来。天主在你中间，你必不动摇，因为至高者已使祂帐幕的场所成圣。因为在你内，天主与祖先所立的盟约和誓愿得到了最荣耀的应验，因为藉着你，主——万军的天主——与我们一同显现了。那不可触摸的荆棘丛，\BibleRef{出 3:2}，预先影射了你那赋予天主尊威的形象，它承载了天主而不被焚毁，天主向先知显现的程度恰如祂愿意被看见的程度。再者，那坚硬粗糙的磐石，\BibleRef{出 17:6}，映照出从你涌向全世界的恩宠与安慰，在旷野中从干渴的侧面涌出丰富的疗愈之泉，供给了昏厥的百姓。还有，那司铎的棍杖，未经栽培却开花结果，\BibleRef{户 17:8}，作为永久司铎职的抵押和定金，为你超自然的生育提供了不容轻视的象征，\BibleRef{希 9:4}。不仅如此？大能者\Person{梅瑟}不是明确宣告过，因着这些关于你的、难解的预象，\BibleRef{出 25:8}，他在山上迟延更久，以便他能学习，啊，圣者，那与你有关系的奥秘吗？因为奉命建造约柜作为这事的标记和相似，他并未疏忽遵命，尽管下山时发生了悲剧；但他按尺寸建造了五肘半的约柜，指定它为法律的容器，并用革鲁宾的翅膀遮盖，最显著地预表了你，天主之母，你无玷地怀了祂，并以不可言喻的方式生下了祂，祂本身可以说就是不朽的本体，而这都在世界五圈半的界限之内。为了你和天主圣言无玷的降生成人，这降生藉着你而发生，为的是那肉躯永不改变、不可分割地与祂永远结合，\BibleRef{希 9:4}。那金罐，作为最确定的预象，保存了其中的玛纳，这玛纳在其他情况下日日改变，却保持不变，历久弥新。先知\Person{厄里亚}\BibleRef{列下 2:11}同样，预知你的贞洁，并藉着圣神效法它，为自己系上了那火焰生命的冠冕，因天主的旨意被判定超越死亡。你也预表了他的继任者\Person{厄里叟}，\BibleRef{德 48:1}，受一位智慧之师的教导，并预先察觉你——尚未诞生的你——的临在，通过某些确信的迹象预示将来要发生的事，向需要的人提供了帮助和医治，这医治具有超自然的能力；现在用新罐盛装疗愈之盐，治愈致命之水，以表明世界要藉着在你内彰显的奥迹被再造；现在用无酵饼，在预表上回应你的生育，未经男人种子玷污，从食物中驱逐死亡的苦涩；然后，通过超越自然的能力，在\Place{约但河}中超越自然元素，从而预先以征兆显示了我们的主下降阴府，以及祂奇妙地将那些被腐败束缚者解救出来。因为万有都屈服于那预表你的神圣肖像之下。\par

十. 但我为何离题，延长我的讲论，任由它被这些各样比喻所驾驭呢？而此时你的事理如同柱子般立于眼前，在其中畅游并欢喜，岂不更好、更有益？因此，辞别历代圣人的属灵叙述和奇事，我转向你——你常被纪念，仿佛把持着这节庆的舵。\par

你有福了，全福者，众所渴望者。主祝福了你的名字，充满天主的恩宠，极其蒙天主悦纳，天主之母，你光照信友。你是那不可界定的者的界定，可谓\OriginalTerm{界定}{circumscription}；那最美丽之花的根\BibleRef{依 40:1}；造物主的母亲；养育者的乳母；那怀抱万有者的四周；那以祂的圣言支撑万有者的支持者\BibleRef{希 1:3}；天主借以在肉身显现的门\BibleRef{则 44:2}；那洁净火炭的火钳\BibleRef{依 6:6}；那包容万有之胸怀的小胸怀；羊毛\BibleRef{民 6:37}，其奥妙无法解释；白冷的水井\BibleRef{撒下 23:17}，达味渴望的生命之泉，从中涌出永生之饮；施恩座\BibleRef{出 35:17}，天主以人形从中向人彰显；那以光为衣者的无瑕长袍。你将那祂所没有的肉身借给一无所缺的天主，为使全能者成为祂乐意成为的。有什么比这更辉煌？有什么比这更崇高？那充满天地、万有都属于祂的\BibleRef{耶 23:24}，竟成为需要你，因为你将祂所没有的肉身借给天主。你以那美丽的身体甲胄披戴了全能者，藉此祂才能被我双眼看见。而我，为能自由地接近瞻仰祂，已领受了那能熄灭恶者一切火箭的\BibleRef{弗 6:16}。万福！万福！天主之母与婢女。万福！万福！万有的大债主向你欠债。我们都欠天主的债，但祂自己却欠你的债。\par

因为那说过：\Quote{孝敬你的父亲和你的母亲，}\BibleRef{出 20:12}的，祂既然愿意藉此证明自己，必定会向那服侍了祂自愿屈尊的诞生者保全那恩宠和祂自己的命令，并会以神圣的尊荣荣耀那一位——祂自己，作为无父者，甚至如同她无丈夫一样，亲自立为母亲。事情必定如此。因为我们将颂歌献给你，至圣可敬的天主居所，并非无益和装饰的话语。同样，对你的属灵赞美也非世俗的轻浮，或虚假谄媚的呼喊——你这受天主赞美者，你给天主哺乳，你以诞生赐予凡人存在的开端——而是清晰明显的真理。但时间会不够我们，世代和后代也不够，向你这位永生之王的母亲\BibleRef{弟前 1:17}献上相称的问候，正如某处那位杰出的先知所说，教导我们你是何等不可测度。天主的殿何其伟大！祂产业的处所何其广阔！伟大，无穷尽，高而不可测量。因为实实在在，这先知的神谕和至真之言，是论及你的尊威；因为你独自被认为配与天主分享天主的事物；你独自在肉身中怀了那位自天主圣父而来、永恒而独生的。凡持守纯正信仰的人，确实如此相信。\par

十一、\newline{}在余下的时间里，我最专注的听众，让我们接过这位老人，这位领受天主者，这位虔诚的导师；他仿佛已从那童贞之海中安全驶抵此地，让我们使他歇息，既满足他的神圣渴望，又将这至为有福的神学传授给我们；我们自己也要继续完成余下的讲论，毫厘不差地指向我们所规定的目标，且在天主全能者的引导下，如此我们便不至于全然无果无益，辜负那对我们所要求的。当预言和司祭职共同被召参加这些神圣礼仪，那一对由天主所拣选的义人——我指的是西默盎和亚纳——身上最清楚地带着两方百姓的表像，已站在这荣耀的童贞宝座旁边，因为那老人代表以色列百姓和如今已衰老的法律，而那位寡妇代表外邦人的教会，她直到此刻仍是寡妇——那老人，作为法律的化身，请求离去；那寡妇，作为教会的化身，献上她喜乐的信仰宣认，并向所有在耶路撒冷期待救赎的人谈论他，正如关于两者的记载都贴切而非凡，与这神圣的庆节完全相合。因为那老人如此精确地知道法律的那道谕令，其中说道：\BibleQuote{你们要听从他，凡不听从他的灵魂，必从民中剪除}，他理当寻求从法律的监护中和平解职；因为当君王在场并向民众说话时，他的一个侍从竟在他对面发言，而民众竟向此人侧耳，这实在是傲慢和僭越。那寡妇既已蒙受超乎限度的恩赐，也理当以节日的曲调向天主谢恩；所以那里发生的事都合乎法律。至于剩下的，我们需要探究：既然预言的预像和象征与这显赫的庆节有某种类比和关系，如已表明的，那么为何说那房屋充满了烟雾？先知并非偶然提及此事，而是意味深长，指的是由天上色辣芬所发出的那\OriginalTerm{至圣者}{Thrice-Holy}的呼声。你若拿起并考察这段叙述之后的内容，就会发现其中的意义，我专注的听众：\BibleQuote{你们听了，却不会明白；看了，却不会察觉。}因此，当愚昧的犹太子孙们看见了那荣耀的奇事——正如达味所歌唱的，主在世上所行的——并看见那从深处与高处相遇的记号，没有分裂，没有混乱；又如依撒意亚先前所宣告的：一位超自然的母亲和一个超理性的孩子；一位地上的母亲和一位天上的儿子；天主对人的本性的一种新的摄取，以及无婚姻的生育；在受造物的循环中有什么比这更荣耀、更值得传颂的呢！然而他们看见了，却如同没有看见；他们闭上眼睛，在赞美上漠不关心。因此他们所夸耀的房屋充满了烟雾。\par

十二、\newline{}此外，除这景象之外，甚至超越这景象，他们听见一位极其正义、极其可信、也值得效法的老人，被圣神感动，是法律的导师，享有司祭的尊荣，因预言之恩赐而显赫，因对基督所怀的希望而延长生命的界限、推迟死亡的债务——当他们看见他，我说，欢喜跳跃，口出吉祥之言，心喜若狂，完全沉浸在神圣而狂喜的激情中，他因敬虔的改变而由人变为天使，因极度的喜乐而高唱感恩之歌，公开宣告：\BibleQuote{为启示外邦人的光明，和祢百姓以色列的荣耀。}即使到那时，他们仍不愿听那放在他们耳中、为天上者自身所尊崇的话；因此他们所夸耀的房屋充满了烟雾。烟雾是愤怒的标志和确凿的证据，正如经上所载：\BibleQuote{祂的怒气上升，烟雾腾起，火从祂面前烧尽；}在另一处：\BibleQuote{在不顺从的民中，火要燃烧。}这火在我们的主所尊崇的福音中清楚地表明了，当祂对犹太人说：\BibleQuote{看，你们的房屋将留给你们荒凉。}又在另一处：\BibleQuote{君王就派军队，消灭那些凶手，焚毁他们的城。}这就是犹太人不信的恶报，他们拒绝向天主圣三献上赞美的归因。因为自从地极被圣化，教会的巍峨房屋因那\OriginalTerm{至圣者}{Thrice Holy}的宣告而充满了主的荣耀，正如大水覆盖海洋，他们便遭遇了先前所宣告的事，预言的开始由其结局得到证实，真理的宣讲者被圣神所感动，如已说的，以实例预示那将要临到他们的可怕毁灭，说道：\BibleQuote{乌齐雅王驾崩那年，我看见主}——乌齐雅无疑作为背逆者，被视为整个背逆群体的代表——他无疑是这群体的头——他也因自己的僭妄而受了惩罚，额上如同铜像般刻着天主的惩罚，以癞病的可憎形象向所有人展示他们那可憎不敬的报应。因此怀有预先知道这些事的天主上智，让感恩的亚纳被接纳，同时让忘恩负义的会堂被排斥。她的名字本身也预先象征教会，那因基督和天主的恩宠在圣洗中成义的教会。因为亚纳按解释是\OriginalTerm{恩宠}{Grace}。\par

XIII. 但在这里，如同在港口一般，我们将那艘带着十字架旗帜的船只停泊，让我们收起演说的风帆，以使它与自身相称。首先，仅以尽可能简短的言词，让我们向大王的城，连同整个教会身体致敬，因为我们在精神上与她们同在，并与圣父和那里备受尊敬的长兄们一同庆节。欢欣吧，大王的城，我们的救恩奥秘在其中得以完成。欢欣吧，你地上的天，\Place{熙雍}，永远忠于主的城。欢欣吧，并照耀吧，\Place{耶路撒冷}，因为你的光明已经来到——永恒的光，永远持久的光，至高的光，非物质的光，与天主圣父同一实质的光，在圣神内的光，圣父也在其中；那光照时代的光；那照明尘世和超尘世事物的光，我们的天主\Person{基督}。欢欣吧，主所神圣和拣选的城。欢乐地庆祝你们的庆节吧，因为它们不会增多以致衰老消逝。欢欣吧，最幸福的城，因为论及你的事是荣耀的；你的司铎要穿上正义，你的圣人们要欢呼，你的穷人将得饱足。欢欣！欢乐吧，\Place{耶路撒冷}，因为主在你中间为王。我说，那位在其单纯而非物质的天主性中，进入了我们的本性，并从童贞女的胎中不可言喻地降生成人的主；那位除了被蛇绊倒的\Person{亚当}的团块之外，别无分享的主。因为主并未抓住天使的种子——我说，那些没有从起初所赋予他们的美好秩序和等级中堕落的天使。祂屈尊就卑地来到我们这里，那位永远与圣父同等的圣言天主。祂来到世间也不是为了恢复，也不会恢复，正如某些邪恶的魔鬼的拥护者所想象的那样，那些曾经从光明中堕落的邪魔；但当万物的创造者和塑造者，正如最神圣的\Person{保禄}所说，抓住了\Person{亚巴郎}的后裔，并通过他抓住了全人类时，祂就永远地、不变地成了人，为的是藉着祂与我们的共融，以及我们与祂的结合，阻止罪恶进入我们中间，其力量逐渐被打破，本身如蜡被熔化，被主来到世间时投在地上的火所熔化。欢欣吧，你，天主教会，你被栽种在全地之上，并与我们一同欢乐。小小羊群，不要害怕敌人的风暴，因为你们的圣父乐意把天国赐给你们，并让你们践踏敌人的颈项。欢欣并欢乐吧，你昔日荒芜、无虔诚后裔，如今却有许多信德之子的女人。欢欣吧，主的子民，你们是特选的种族，王家的司铎，圣洁的国民，属于主的民族——要宣扬那召叫你们出离黑暗进入奇妙光明的祂的赞美，并为祂的仁慈而光荣祂。\par

XIV. 万福，永远的童贞天主之母，我们不断的喜乐，因为我再次回到你面前。你是我们庆节的开端、中间和终结；属于天国的重价珍珠；每一祭品的脂油，生命之粮的活祭台。万福，天主之爱的宝藏。万福，圣子爱人的泉源。万福，圣神的遮荫山。你闪耀着，甜蜜的赐恩之母，如同太阳的光芒；你闪耀着最炽热爱德的不可承受的火焰，在末后产生了那在起初以前就由你孕育的，显明了那隐藏而不可言喻的奥秘，圣父的不可见之子——和平之王，祂以奇妙的方式显示自己比一切微小更微小。因此，我们恳求你，妇女中最卓越者，你以母性尊荣的自信心夸耀，请不断纪念我们。天主之母，请纪念我们，我说，以你为夸耀，并以庄严的颂歌庆祝那永存不朽的纪念。还有，可敬可颂的\Person{西默盎}，我们神圣宗教的最初接待者，信众复活的导师，请作我们的主保和代祷者，向那位你配得抱在怀中的救主天主。我们与你一同向那掌管生死的\Person{基督}歌唱赞美，说：你是从真光而来的真光；从真天主而来的真天主；唯一的主，在祢取得人性之前；然而在祢取得人性之后，那同一的主，永远当受朝拜；凭自身是天主，而非因恩宠，但为我们也是完全的人；在本性上绝对君权的主，但为我们和我们的救恩也取了仆人的形态，然而无玷无染。因为祢是不朽的，来释放朽坏，为使万物成为不朽。因为光荣、权能、伟大、威严都属于祢，与圣父和圣神，至于永远。阿们。\par

\SourceRecord{Translated by William R. Clark.；Ante-Nicene Fathers；Vol. 6.；Edited by Alexander Roberts, James Donaldson, and A. Cleveland Coxe.；Buffalo, NY: Christian Literature Publishing Co.,；1886.；<http://www.newadvent.org/fathers/0627.htm>.}
